
\documentclass[tikz]{standalone}
\definecolor{bg}{HTML}{171D2B}
\definecolor{ak}{HTML}{F7B538}
\begin{document}
\pagecolor{bg}
\begin{tikzpicture}[line cap=round, text=white]
  % jam pertama (kiri): 3.00 -> 90 derajat
  \begin{scope}[shift={(0,0)}]
    \draw[white, line width=2.5pt] (0,0) circle (2.2);
    \foreach \a in {0,30,...,330}
      \draw[white!60, line width=1pt] (\a:1.95) -- (\a:2.2);
    \foreach \n/\a in {12/90,3/0,6/270,9/180}
      \node[white] at (\a:1.6) {\n};
    \draw[ak, line width=3pt] (0,0) -- (0,1.4);        % jarum menit di angka 12
    \draw[white, line width=4pt] (0,0) -- (1.2,0);     % jarum jam di angka 3
    \fill[ak] (0,0) circle (0.09);
    \node[white] at (0,-2.75) {Pukul 03.00 $\rightarrow$ $90^\circ$};
  \end{scope}
  % jam kedua (kanan): 5.00 -> 150 derajat, busur ditandai
  \begin{scope}[shift={(7.2,0)}]
    \draw[white, line width=2.5pt] (0,0) circle (2.2);
    \foreach \a in {0,30,...,330}
      \draw[white!60, line width=1pt] (\a:1.95) -- (\a:2.2);
    \foreach \n/\a in {12/90,3/0,6/270,9/180}
      \node[white] at (\a:1.6) {\n};
    \draw[ak, line width=3pt] (0,0) -- (90:1.4);       % menit di 12
    \draw[white, line width=4pt] (0,0) -- (-60:1.2);   % jam di 5
    \draw[ak, line width=1.5pt] (0.75,0) arc[start angle=0, end angle=-60, radius=0.75];
    \node[ak] at (-25:1.25) {$150^\circ$};
    \node[white] at (0,-2.75) {Pukul 05.00 $\rightarrow$ $150^\circ$};
  \end{scope}
\end{tikzpicture}
\end{document}
